% ============================================================
%  Bonus, Segment 3B — Functional Strumming Patterns / 功能性扫弦模式
%  原书页码 82-92
% ============================================================
\chapter{Functional Strumming Patterns / 功能性扫弦模式}
\label{ch:strumming}

% ---- 此文件保留了原有已翻译内容 ----

\section{Overview / 概述}

\eng{Strumming patterns are the foundation of rhythm guitar in math rock.
Unlike traditional rock or pop, math rock often employs odd time signatures
and syncopated rhythms.}

\chn{扫弦节奏是数学摇滚节奏吉他的基础。
与传统摇滚或流行不同，数学摇滚常用奇数拍号和切分节奏。}

\enandcn
  {Let's start with common time signatures.}
  {从最常用的拍号开始。}

\paragraph{核心术语}
\begin{itemize}[leftmargin=*]
  \item \term{奇数拍号}{Odd Time Signature}
  \item \term{切分音}{Syncopation}
  \item \term{连音符}{Tuplet}
\end{itemize}

\begin{notebox}[提示]
  数学摇滚扫弦常强调反拍。练习时先用节拍器慢速感受重音移位。
\end{notebox}

\section{3/4 Time / 四三拍}

\begin{examplebox}[3/4 扫弦模式]
  3/4 拍中常见的扫弦模式：
\end{examplebox}

\begin{verse}
  \small\ttfamily
  e|-----------------|-----------------|\\
  B|-----------------|-----------------|\\
  G|-----------------|-----------------|\\
  D|-------2---------|-------2---------|\\
  A|-----3---3-------|-----3---3-------|\\
  E|---3-------3-----|---3-------3-----|\\
\end{verse}

\eng{This pattern uses palm muting on the downbeats with open strings ringing.}
\chn{这一模式使用下拍制音，空弦音延续。}

\begin{exercise}[基础扫弦]
  以 BPM = 80 练习上述模式：
  \begin{enumerate}
    \item 保持手腕放松
    \item 强调第 2 拍反拍
    \item 先慢后快，逐步加速
  \end{enumerate}
\end{exercise}

\section{4/4 Time / 四四拍}

\eng{In 4/4, math rock strumming creates tension
by emphasizing the \& of 2 and \& of 4.}
\chn{在 4/4 拍中，数学摇滚扫弦通过强调第 2、4 拍反拍制造张力。}

% ---- 后续翻译内容 ----
% 原书还包括：
%   - 4/4 Time: 13 种扫弦模式 (5A-5M)
%   - 3/4 Time: 6 种扫弦模式 (6A-6F)
%   - 5/8 Time: 8 种扫弦模式 (7A-7H)
%   - 附录：时值划分速查表
